\section{Related Work}
\label{sec:relatedwork}

\noindent \textbf{Embodied AI platforms.} Various Embodied AI platforms have been developed over the past several years~\cite{kolve2017ai2,savva2019habitat,Shen2020iGibsonAS,xiang2020sapien,Gan2020ThreeDWorldAP,house3d}. These platforms target different design goals. AI2-THOR~\cite{kolve2017ai2} and its variants (ManipulaTHOR~\cite{manipulathor} and RoboTHOR~\cite{Deitke2020RoboTHORAO}) are built in the Unity game engine and focus on agent-object interactions, object state changes, and accurate physics simulation. Unlike AI2-THOR, Habitat~\cite{savva2019habitat} provides scenes constructed from 3D scans of houses, however, objects and scenes are not interactable. A more recent version, Habitat 2.0~\cite{szot2021habitat}, introduces object interactions at the expense of being limited to one floorplan and synthetic scenes. iGibson~\cite{Shen2020iGibsonAS} includes photo-realistic scenes, but with limited interactions such as pushing. iGibson 2.0~\cite{Li2021iGibson2O} extends iGibson by focusing on household tasks and object state changes in synthetic scenes and includes a virtual reality interface. ThreeDWorld~\cite{Gan2020ThreeDWorldAP} targets high-fidelity physics simulation such as liquid and deformable object simulation. VirtualHome~\cite{Puig2018VirtualHomeSH} is designed for simulating human activities via programs. RLBench~\cite{james2020rlbench}, RoboSuite~\cite{zhu2020robosuite} and Sapien~\cite{xiang2020sapien} target fine-grained manipulation. The main advantage of \env is that we can generate a diverse set of \emph{interactive} scenes procedurally, enabling studies of data augmentation and large-scale training in the context of Embodied AI. 
 
\noindent \textbf{Large-scale datasets.} Large-scale datasets have resulted in major breakthroughs in different domains such as image classification~\cite{imagenet,Kuznetsova2020TheOI}, vision and language~\cite{Changpinyo2021Conceptual1P,Thomee2016YFCC100MTN}, 3D understanding~\cite{shapenet2015,Xiang2016ObjectNet3DAL}, autonomous driving~\cite{nuscenes,sun2020scalability}, and robotic object manipulation~\cite{pinto2016supersizing,mu2021maniskill}. However, there are not many interactive large-scale datasets for Embodied AI research. \env includes interactive houses generated procedurally. Hence, there are an arbitrarily large number of scenes in the framework. The closest works to ours are~\cite{Ramakrishnan2021HabitatMatterport3D,Petrenko2021MegaverseSE,li2021openrooms}. HM3D~\cite{Ramakrishnan2021HabitatMatterport3D} is a recent framework that includes 1,000 scenes generated using 3D scans of real environments. \env has a number of key distinctions: (1) unlike HM3D which includes static scenes, the scenes in \env are interactive i.e., objects can move and change state, the lighting and texture of objects can change, and a physics engine determines the future states of the scenes; (2) it is challenging to scale up HM3D as it requires scanning a house and cleaning up the data, while we can procedurally generate more houses; (3) HM3D can be used only for navigation tasks (as there is no physics simulation and object interaction), while \env can be used for tasks other than navigation. OpenRooms~\cite{li2021openrooms} is similar to HM3D in terms of the source of the data (3D scans) and dataset size. However, OpenRooms is interactive. OpenRooms is also confined to the set of scanned houses, and it takes a significant amount of time to annotate a new scene (e.g., labeling materials for one object takes 1 minute), while \env does not suffer from these issues. Megaverse~\cite{Petrenko2021MegaverseSE} is another large-scale Embodied AI platform that includes procedurally generated environments. Although it is impressive in terms of simulation speed, it includes only game-like environments with a simplified appearance. In contrast, \env mimics real-world houses in terms of the complexity of appearance, physics, and object interactions. 

\noindent \textbf{Scene generation.} Indoor scene synthesis has been studied extensively in computer vision and graphics communities. \cite{Chang2015TextT3,Chang2014LearningSK,chang2014interactive} address generating 3D scenes from text descriptions. \cite{wu2019data,Hu2020Graph2PlanLF,Nauata2021HouseGANGA} learn to generate house floorplans. \cite{Ritchie2019FastAF,Zhang2020DeepGM,Chaudhuri2019LearningGM,Wang2021SceneFormerIS,keshavarzi2020scenegen,li2019grains} use generative models for indoor scene generation. \cite{zhou2019scenegraphnet,Savva2017SceneSuggestC3} propose potential objects for a query location in an indoor scene. Others have used procedural generation \cite{khalifa2020pcgrl, earle2021learning} and unsupervised learning \cite{dennis2020emergent} to synthesize grid-world environments for AI. \env is specifically designed for Embodied AI research in the sense that (1) all scenes are interactive and physics-enabled, and the placement of objects respects the physics of the world (e.g., there are no two objects that clip through each other), (2) there are various forms of scene augmentation such as randomization of object placements while following certain commonsense rules, variation in the appearance of objects and structures, and variation in lighting. %\todo{use better things to differentiate ourselves}


